\begin{onehalfspace}

\section{Il Problema della Fattorizzazione Intera}

La \textbf{fattorizzazione intera} consiste nel trovare i fattori primi di un numero intero $N$. Per numeri piccoli il problema è banale, ma la sua difficoltà cresce in modo impressionante all'aumentare di $N$: i migliori algoritmi classici noti, come il \textit{General Number Field Sieve} (\ac{GNFS}), richiedono un tempo sub-esponenziale ma comunque super-polinomiale:\footnote{Il \ac{GNFS} detiene il record mondiale per la fattorizzazione classica: nel 2020, un consorzio internazionale ha fattorizzato RSA-250 (829 bit, 250 cifre decimali) \cite{rsa250_2020}, impiegando l'equivalente di circa 2700 anni-core. In precedenza, nel 2009, era stato fattorizzato RSA-768 (768 bit). I numeri RSA usati in produzione sono di 2048 o 4096 bit, rendendo la fattorizzazione classica del tutto impraticabile con qualunque tecnologia convenzionale prevedibile.}
\begin{equation}
    T_{\text{classico}}(N) = O\!\left(\exp\!\left(c(\log N)^{1/3}(\log\log N)^{2/3}\right)\right)
\end{equation}

Questo problema è classificato in NP ma non è noto se appartenga a P: non esiste alcun algoritmo classico polinomiale per risolverlo, e si congettura fortemente che non esista.

\subsection{Rilevanza Crittografica: RSA}

Il sistema crittografico \textbf{RSA} (Rivest-Shamir-Adleman, 1977) \cite{rsa1978} è il più diffuso algoritmo a chiave pubblica, alla base di HTTPS, TLS, SSH e di gran parte delle comunicazioni digitali sicure. La sua sicurezza si fonda interamente sulla difficoltà computazionale della fattorizzazione:

\begin{itemize}
    \item Si scelgono due grandi primi $p$ e $q$ (tipicamente 1024--4096 bit ciascuno);
    \item Si calcola il modulo $N = p \cdot q$ (la \textit{chiave pubblica} include $N$ ed un esponente $e$);
    \item La \textit{chiave privata} è determinata da $d$ tale che $ed \equiv 1 \pmod{\lambda(N)}$, dove $\lambda$ è la funzione di Carmichael;
    \item Decifrare un messaggio senza conoscere $p$ e $q$ equivale, in pratica, a fattorizzare $N$.
\end{itemize}

Fattorizzare un numero RSA a 2048 bit richiederebbe, con gli algoritmi classici attuali, un tempo stimato superiore all'età dell'universo. Nel 1994, Peter Shor\footnote{\textbf{Peter Williston Shor} (nato nel 1959 a New York) è professore di matematica applicata al \ac{MIT}. Al momento della scoperta dell'algoritmo lavorava presso i laboratori di ricerca AT\&T Bell Labs. Il suo algoritmo è considerato la dimostrazione più impattante del potere del calcolo quantistico.} dimostrò che un computer quantistico può eseguire questa fattorizzazione in tempo \textit{polinomiale}, rendendo RSA teoricamente vulnerabile \cite{shor1994}.

\section{Riduzione alla Ricerca del Periodo}

L'intuizione geniale di Shor consiste nel ricondurre la fattorizzazione a un problema di \textbf{period finding}, trattabile quantisticamente in modo efficiente.

\begin{description}
    \item[Osservazione chiave:] Dato $N$ e un intero $a$ con $\gcd(a, N) = 1$ (coprimo con $N$), la funzione
    \begin{equation}
        f(x) = a^x \bmod N
    \end{equation}
    è \textbf{periodica} con periodo $r$ (detto \textit{ordine} di $a$ modulo $N$): $f(x+r) = f(x)$ per ogni $x$.
\end{description}

Una volta trovato il periodo $r$, se $r$ è pari si può calcolare:
\begin{equation}
    \gcd\!\left(a^{r/2} - 1, N\right) \quad \text{e} \quad \gcd\!\left(a^{r/2} + 1, N\right)
\end{equation}

Con alta probabilità, almeno uno di questi due valori è un fattore non banale di $N$. La prova utilizza proprietà della teoria dei numeri: in particolare, vale $(a^{r/2}-1)(a^{r/2}+1) = a^r - 1 \equiv 0 \pmod{N}$.

\section{La Quantum Fourier Transform}

Il componente quantistico fondamentale dell'algoritmo di Shor è la \textbf{Quantum Fourier Transform} (QFT), l'analogo quantistico della Trasformata Discreta di Fourier (DFT).

\subsection{Definizione}

La QFT su $\mathbb{Z}_{2^n}$ agisce su uno stato base come:
\begin{equation}
    \text{QFT}\ket{j} = \frac{1}{\sqrt{2^n}} \sum_{k=0}^{2^n-1} e^{2\pi i jk / 2^n} \ket{k}
\end{equation}

L'azione su uno stato generico $\ket{\psi} = \sum_j \alpha_j \ket{j}$ produce:
\begin{equation}
    \text{QFT}\ket{\psi} = \sum_k \hat{\alpha}_k \ket{k}
\end{equation}

dove $\hat{\alpha}_k = \frac{1}{\sqrt{2^n}}\sum_j \alpha_j e^{2\pi i jk/2^n}$ sono i coefficienti della DFT.

\subsection{Implementazione Efficiente}

La DFT classica su $N = 2^n$ elementi richiede $O(N \log N)$ operazioni (con il Fast Fourier Transform). La QFT, invece, richiede solo $O(n^2) = O((\log N)^2)$ porte quantistiche, ottenendo uno speedup esponenziale.

Il circuito della QFT è costruito ricorsivamente tramite porte di Hadamard e porte di fase controllate $CR_k$ dove:
\begin{equation}
    CR_k = \begin{pmatrix} 1&0&0&0\\0&1&0&0\\0&0&1&0\\0&0&0&e^{2\pi i/2^k} \end{pmatrix}
\end{equation}

\subsection{Phase Estimation}

La QFT è impiegata, all'interno dell'algoritmo di Shor, attraverso la subroutine di \textbf{Phase Estimation} (QPE). Data una porta $U$ con autovettore $\ket{u}$ e autovalore $e^{2\pi i \theta}$, la QPE stima $\theta$ con precisione arbitraria in tempo $O(n^2)$. Nell'algoritmo di Shor, $U$ è l'operatore di moltiplicazione modulare $U\ket{x} = \ket{ax \bmod N}$, e il suo spettro di fase contiene l'informazione sul periodo $r$.

\section{L'Algoritmo di Shor: Schema Completo}

\begin{algorithm}[H]
\caption{Algoritmo di Shor per la fattorizzazione di $N$}
\begin{algorithmic}[1]
\Require Numero composto $N$, numero di qubit $n = 2\lceil\log_2 N\rceil$
\Ensure Fattore non banale di $N$, oppure \textit{fallimento}
\State Scegliere $a \in \{2, \ldots, N-1\}$ casuale
\State Se $g = \gcd(a,N) > 1$ allora \Return $g$ \Comment{Fortuna: $a$ già condivide un fattore}
\State Inizializzare $\ket{0}^{\otimes n}\ket{1}$
\State Applicare $H^{\otimes n}$ al primo registro: $\frac{1}{\sqrt{2^n}}\sum_{x=0}^{2^n-1}\ket{x}\ket{1}$
\State Applicare l'oracolo di moltiplicazione modulare: $\frac{1}{\sqrt{2^n}}\sum_{x=0}^{2^n-1}\ket{x}\ket{a^x \bmod N}$
\State Misurare il secondo registro: il primo collassa in sovrapposizione degli $x$ con stesso residuo
\State Applicare la QFT inversa al primo registro
\State Misurare il primo registro: ottenere una stima di $s/r$
\State Usare l'algoritmo delle frazioni continue per determinare $r$ da $s/r$
\If{$r$ è dispari o $a^{r/2} \equiv -1 \pmod{N}$} \Return \textit{Fallimento}
\Else
    \State Calcolare $p = \gcd(a^{r/2}-1, N)$ e $q = \gcd(a^{r/2}+1, N)$
    \If{$p$ o $q$ è un fattore non banale di $N$} \Return $p$ o $q$
    \EndIf
\EndIf
\end{algorithmic}
\end{algorithm}

\section{Analisi della Complessità}

La complessità dell'algoritmo di Shor è \textbf{polinomiale} nel numero di bit di $N$:
\begin{itemize}
    \item Il circuito quantistico richiede $O((\log N)^2 \log\log N)$ porte elementari;
    \item Il numero di qubit necessari è $O(\log N)$;
    \item Il numero atteso di iterazioni (con basi $a$ diverse) prima del successo è $O(1)$.
\end{itemize}

Complessivamente: $T_{\text{Shor}}(N) = O\!\left((\log N)^3\right)$, a fronte del $O(\exp((\log N)^{1/3}))$ del miglior algoritmo classico. Si tratta di uno \textbf{speedup esponenziale}.

\section{Impatto sulla Crittografia e Prospettive Post-Quantum}

L'algoritmo di Shor rende teoricamente vulnerabili tutti i sistemi crittografici basati sulla difficoltà della fattorizzazione (RSA) o del logaritmo discreto (Diffie-Hellman, DSA, ECDSA). La \textit{Post-Quantum Cryptography} (PQC) è il campo emergente che sviluppa algoritmi crittografici resistenti agli attacchi quantistici.

Il processo di standardizzazione \ac{PQC} del \ac{NIST}, avviato nel 2016 con una call pubblica per la raccolta di proposte algoritmiche, si è concluso con la pubblicazione formale dei primi standard federali nell'agosto 2024 \cite{nist2024pqc}:\footnote{Dopo otto anni di valutazione e quattro round di selezione pubblica con la partecipazione di crittografi di tutto il mondo, il NIST ha pubblicato nel 2024 i FIPS 203, 204 e 205. Un quarto standard, FIPS 206 (FN-DSA, basato su FALCON), è in fase di finalizzazione. I nuovi nomi riflettono le famiglie algoritmiche di appartenenza: ML-KEM (Module-Lattice Key Encapsulation), ML-DSA (Module-Lattice Digital Signature), SLH-DSA (Stateless Hash-based Digital Signature).}
\begin{itemize}
    \item \textbf{FIPS 203 -- ML-KEM} (basato su CRYSTALS-Kyber): standard per l'incapsulamento di chiave, basato su reticoli modulari;
    \item \textbf{FIPS 204 -- ML-DSA} (basato su CRYSTALS-Dilithium): standard per la firma digitale su reticoli;
    \item \textbf{FIPS 205 -- SLH-DSA} (basato su SPHINCS+): standard per la firma digitale basata su funzioni hash;
    \item \textbf{FIPS 206 -- FN-DSA} (basato su FALCON): standard in fase di finalizzazione per la firma digitale su reticoli NTRU.
\end{itemize}

Questi sistemi si basano su problemi matematici (reticoli, hash, codici) per i quali non sono noti algoritmi quantistici efficienti, garantendo sicurezza anche nell'era post-quantum. La transizione verso questi standard è urgente anche per il cosiddetto fenomeno \textit{harvest now, decrypt later}: attori malevoli raccolgono oggi comunicazioni cifrate con RSA, con l'intenzione di decifrarle quando saranno disponibili computer quantistici fault-tolerant.

\section{Sviluppi Recenti e Stato dell'Arte}

\subsection{Implementazioni su Hardware Reale}

Nonostante la solidità teorica, l'esecuzione di Shor su hardware reale rimane limitata a istanze molto piccole a causa dei requisiti di profondità circuitale. I principali risultati sperimentali sono:

\begin{itemize}
    \item \textbf{2001 -- NMR (Vandersypen et al.)}: prima dimostrazione sperimentale di Shor su un sistema a 7 spin nucleari; fattorizzazione di $N = 15$ in $3 \times 5$ \cite{vandersypen2001}.
    \item \textbf{2016 -- Ioni intrappolati (Monz et al.)}: implementazione scalabile su 5 qubit ionici; fattorizzazione di $N = 15$ con un'architettura modulare \cite{monz2016}.
    \item \textbf{2021 -- Qubit superconduttori IBM}: dimostrazione della fattorizzazione di $N = 21$ su processori IBM Quantum utilizzando solo 5 qubit \cite{skosana2021}, attualmente il risultato più avanzato su hardware superconduttore di uso generale;
    \item \textbf{Fotonica}: esperimenti su piattaforme fotonice hanno raggiunto la fattorizzazione di $N = 35$, sfruttando la maggiore coerenza dei qubit ottici, sebbene con circuiti fortemente semplificati rispetto all'algoritmo generale.
\end{itemize}

Il record di $N = 21$ evidenzia quanto la distanza tra le istanze dimostrabili e quelle crittograficamente rilevanti (RSA-2048, con $N$ a 2048 bit) sia ancora enorme: fattorizzare un numero a $L$ bit richiede circa $5L+1$ qubit logici e $72L^3$ gate \cite{shor_challenges2025}, valori del tutto inaccessibili ai processori NISQ attuali. Una analisi aggiornata delle sfide pratiche nell'esecuzione di Shor su piattaforme esistenti è discussa in \cite{shor_challenges2025}.

\subsection{Miglioramenti Algoritmici: Regev 2023}

Un risultato teorico rilevante è stato pubblicato nel 2023 da Oded Regev \cite{regev2023}, che ha proposto una variante migliorata dell'algoritmo di Shor richiedente solo $O(n^{3/2})$ gate (rispetto agli $O(n^2)$ della versione originale), al costo di un overhead classico maggiore. Sebbene non ancora implementata su hardware reale, questa riduzione della profondità circuitale è promettente per le architetture NISQ, dove la profondità è il principale collo di bottiglia.

\subsection{Roadmap IBM verso il Quantum Fault-Tolerant}

IBM ha delineato una roadmap concreta verso il calcolo quantistico fault-tolerant \cite{ibm_roadmap2025}:
\begin{itemize}
    \item \textbf{2025}: processori Loon e Nighthawk (120 qubit con 218 coupler sintonizzabili), con dimostrazione dei componenti chiave per la fault-tolerance tramite codici qLDPC;
    \item \textbf{2026}: target di vantaggio quantistico su problemi pratici;
    \item \textbf{2029}: sistema Quantum Starling con 200 qubit logici e oltre $10^8$ operazioni quantistiche.
\end{itemize}

Questi traguardi definiscono l'orizzonte temporale entro cui l'algoritmo di Shor potrà diventare una minaccia concreta per la crittografia RSA, motivando ulteriormente la transizione urgente verso gli standard PQC discussi nella sezione precedente.

Le sfide hardware principali per l'esecuzione di Shor su istanze crittograficamente rilevanti sono analizzate in dettaglio nel Capitolo~\ref{chap:rumore}, mentre l'implementazione pratica con Qiskit e le tecniche di ottimizzazione basate su intelligenza artificiale sono presentate nel Capitolo~\ref{chap:implementazioni}.

\section{Contesto: Confronto con l'Algoritmo di Grover}

Per collocare correttamente l'algoritmo di Shor nel panorama dell'informatica quantistica, è utile confrontarlo con l'altro grande risultato algoritmico degli anni Novanta: l'algoritmo di Grover per la ricerca non strutturata \cite{grover1996}.

\subsection{Il Problema di Grover}

Dato un insieme di $N$ elementi e una funzione oracolo $f:\{0,\ldots,N-1\}\to\{0,1\}$, l'obiettivo è trovare $x^*$ tale che $f(x^*)=1$. Classicamente, nel caso peggiore, sono necessarie $O(N)$ interrogazioni. Grover (1996) dimostrò che un algoritmo quantistico, sfruttando il meccanismo di \textit{amplitude amplification}, può trovare la soluzione in sole $O(\sqrt{N})$ interrogazioni -- uno \textbf{speedup quadratico} che è stato dimostrato ottimale: nessun algoritmo quantistico può fare meglio \cite{bennett1997,zalka1999}.

\subsection{Differenze Fondamentali}

I due algoritmi rappresentano due regimi distinti di vantaggio quantistico:

\begin{table}[H]
\centering
\caption{Confronto tra l'algoritmo di Shor e l'algoritmo di Grover}
\begin{tabular}{@{}lcc@{}}
\toprule
\textbf{Caratteristica} & \textbf{Shor} & \textbf{Grover} \\
\midrule
Problema target          & Fattorizzazione / periodo       & Ricerca non strutturata \\
Speedup                  & Esponenziale ($O(n^3)$ vs sub-exp) & Quadratico ($O(\sqrt{N})$ vs $O(N)$) \\
Tecnica quantistica core & QFT + Phase Estimation          & Amplitude Amplification \\
Ottimalità               & Non dimostrata in assoluto      & Dimostrata ottimale \\
Impatto crittografico    & Rompe RSA, DH, ECDSA            & Dimezza sicurezza AES/SHA \\
Requisiti hardware       & $O(\log N)$ qubit + alta profondità & $O(\log N)$ qubit + bassa profondità \\
\bottomrule
\end{tabular}
\end{table}

L'impatto sulla crittografia simmetrica di Grover è significativo ma gestibile: raddoppiare la lunghezza delle chiavi (da AES-128 a AES-256) è sufficiente a ripristinare il livello di sicurezza. Al contrario, non esiste un analogo rimedio per RSA: l'unica risposta è la migrazione verso gli standard post-quantum discussi nella sezione precedente. È per questa ragione che l'algoritmo di Shor rappresenta la minaccia quantistica più urgente e motivante per la ricerca sul calcolo quantistico applicato.

\end{onehalfspace}
